\paragraph{Buyer valuations and demand.}
Section~\ref{ss:model_setup} writes the valuation of an EPC-sensitive buyer in what we call \emph{bonus form}: she adds $\beta_1$ times her belief that the dwelling is efficient. This form records the difference in energy value between efficient and inefficient dwellings, $\beta_1$, but never the level of the energy cost of either, and it has no place for a buyer who forms no belief at all, because every sensitive buyer enters through some number $\mathbb E[\theta\mid\cdot\,]\in[0,1]$. The ceiling documented in Section~\ref{ss:att_motivation} is thus a property of the representation as much as of the economics. To ask what a buyer pays for a listing whose energy performance she does not consider, we rewrite the valuation in \emph{cost form}. The rewriting changes nothing in the model of Section~\ref{ss:model_setup} (Lemma~\ref{lem:cost_form}); it changes what the model can express. Appendix~\ref{app:attention} explains why the normalisation implicit in the bonus form becomes binding once attention depends on what a listing shows.
